\chapter{Wound care}
\index{Wound care}
\section*{Definition and Overview}
Wound care is the assessment and treatment of damaged skin or deeper tissue, with the aims of controlling bleeding, preventing infection, preserving function, and supporting healing. Care differs for a clean acute wound, a burn, a bite, a surgical wound, and a chronic ulcer.
\section*{Clinical Features}
Important findings include wound size and depth, contamination, bleeding, pain, tissue colour, sensation, movement, pulses, foreign material, and signs of infection. Document changes over time, especially in chronic wounds.
\section*{When to Seek Medical Care}
Seek urgent care for uncontrolled bleeding, deep or gaping wounds, bites, punctures, wounds involving an eye, face, hand, joint, tendon, nerve or bone, severe burns, numbness, loss of movement, or any wound with fever or rapidly spreading redness. People with diabetes, poor circulation, immune suppression, or recurrent wounds need a lower threshold for assessment.
\section*{Management}
For a minor wound, wash hands, apply direct pressure, rinse with clean running water, remove loose debris, and cover with a sterile or clean dressing. Change the dressing when wet, dirty, or as advised. Avoid picking scabs and avoid putting alcohol, peroxide, or household remedies into the wound. Professional care may include closure, debridement, drainage, antibiotics for selected infections, tetanus vaccination, pressure relief, vascular assessment, and specialist wound review.
\section*{Complications}
Watch for increasing pain, warmth, swelling, pus, spreading redness, fever, wound breakdown, black or dead tissue, red streaks, or loss of function. Complications include cellulitis, abscess, sepsis, scarring, delayed healing, and tissue or limb loss.
\section*{Prognosis and Outcomes}
Most minor wounds heal with simple care. Chronic or poorly perfused wounds may take weeks or months and require coordinated treatment of the underlying cause.
\section*{Prevention and Self-Care}
Keep tetanus vaccination current, use protective equipment, inspect feet if diabetes or reduced sensation is present, reduce pressure and friction, stop smoking, maintain nutrition, and attend planned reviews.
\section*{Sources and Further Reading}
\begin{itemize}
    \item NHS cuts and grazes: \url{https://www.nhs.uk/conditions/cuts-and-grazes/}
\end{itemize}
